\section{Модель контактной динамики и решение ограничений}

Алгоритмы обнаружения столкновений (раздел~2.2) предоставляют геометрическую информацию о контакте.
На её основе задача динамики контактирующих тел сводится к задаче поиска сил или импульсов,
которые предотвращают взаимопроникновение и моделируют трение.
В данном разделе изложена импульсная формулировка этой задачи как линейной задачи комплементарности (LCP) и описан итерационный метод её решения.

\subsection{Скоростные ограничения}

Вместо того чтобы работать с силами (что требует малых шагов интегрирования),
используют \emph{импульсный подход}: на каждом шаге симуляции к телам применяются мгновенные импульсы,
изменяющие скорости~\cite{catto2005iterative}.

Для контакта двух тел $A$ и $B$ с нормалью $\hat{\mathbf{n}}$ (направленной от $B$ к $A$) и точками контакта $\mathbf{p}_A$, $\mathbf{p}_B$ определим
относительную скорость в контактной точке:
\begin{equation}
  v_{\text{rel}} = \hat{\mathbf{n}} \cdot \left(
    \mathbf{v}_A + \boldsymbol{\omega}_A \times \mathbf{r}_A - \mathbf{v}_B - \boldsymbol{\omega}_B \times \mathbf{r}_B
  \right),
  \label{eq:vrel}
\end{equation}
где $\mathbf{r}_A = \mathbf{p} - \mathbf{c}_A$ и $\mathbf{r}_B = \mathbf{p} - \mathbf{c}_B$ — радиус-векторы точки контакта относительно центров масс.

Условие непроникновения в скоростной форме:
\begin{equation}
  v_{\text{rel}} + e\, v_{\text{rel}}^{-} \geq 0,
  \label{eq:non-penetration}
\end{equation}
где $v_{\text{rel}}^{-}$ — относительная скорость до применения импульса, $e \in [0, 1]$ — коэффициент упругости.

\subsection{Якобиан ограничения}

В общей форме ограничение $k$ можно записать как:
\begin{equation}
  C_k(\mathbf{q}) = 0 \quad \Longrightarrow \quad \mathbf{J}_k\,\mathbf{v} + b_k = 0,
  \label{eq:constraint}
\end{equation}
где $\mathbf{J}_k \in \mathbb{R}^{1 \times 12}$ — якобиан ограничения, $\mathbf{v} = (\mathbf{v}_A, \boldsymbol{\omega}_A, \mathbf{v}_B, \boldsymbol{\omega}_B)^T$, $b_k$ — поправочный член (Baumgarte stabilization).

Для нормального контактного ограничения якобиан имеет вид:
\begin{equation}
  \mathbf{J}_n = \left[\hat{\mathbf{n}}^T,\; (\mathbf{r}_A \times \hat{\mathbf{n}})^T,\; -\hat{\mathbf{n}}^T,\; -(\mathbf{r}_B \times \hat{\mathbf{n}})^T\right].
  \label{eq:jacobian}
\end{equation}

Для касательных ограничений (трение) используются два якобиана $\mathbf{J}_{t1}$, $\mathbf{J}_{t2}$ по касательным $\hat{\mathbf{t}}_1$, $\hat{\mathbf{t}}_2$ в контактной плоскости.

\subsection{Решение системы ограничений. Обобщённый импульс}

Пусть имеется $m$ ограничений.
Суммарный импульс, действующий на тело $A$ от всех ограничений:
\begin{equation}
  \Delta\mathbf{v}_A = \mathbf{M}_A^{-1}\sum_{k=1}^m \mathbf{J}_{kA}^T \lambda_k,
  \label{eq:delta-v}
\end{equation}
где $\lambda_k$ — \emph{множитель Лагранжа} ограничения $k$ (для нормального контакта — нормальный импульс),
$\mathbf{M}_A^{-1} = \text{diag}(m_A^{-1}\mathbf{E}, \mathbf{I}_A^{-1})$ — обобщённая обратная матрица массы.

Подставляя~\eqref{eq:delta-v} в условие~\eqref{eq:constraint}, получаем систему линейных уравнений:
\begin{equation}
  \mathbf{A}\,\boldsymbol{\lambda} = \mathbf{b},
  \label{eq:lcp-linear}
\end{equation}
где $A_{ij} = \mathbf{J}_i\,\mathbf{M}^{-1}\mathbf{J}_j^T$ — \emph{матрица Делассюса}.

\subsection{LCP-формулировка}

Для контактных ограничений задача обладает условиями комплементарности:
нормальный импульс $\lambda_n$ не может быть отрицательным (тела не прилипают),
и либо нет проникновения, либо нет импульса:
\begin{equation}
  \lambda_n \geq 0, \quad v_{\text{rel}} \geq 0, \quad \lambda_n \cdot v_{\text{rel}} = 0.
  \label{eq:complementarity}
\end{equation}

Условие трения Кулона ограничивает касательный импульс:
\begin{equation}
  \|\boldsymbol{\lambda}_t\| \leq \mu\,\lambda_n,
  \label{eq:coulomb}
\end{equation}
где $\mu$ — коэффициент трения.

Система~\eqref{eq:lcp-linear} с условиями~\eqref{eq:complementarity} является \emph{линейной задачей комплементарности} (LCP).
Для дискретизированного кулонова трения с конечным числом касательных направлений LCP допускает итерационное решение методом Projected Gauss-Seidel.

\subsection{Решатель PGS (Projected Gauss-Seidel)}

PGS~\cite{catto2005iterative} — итерационный решатель, обрабатывающий ограничения по одному, применяя проекцию для учёта неравенств.

На каждой итерации для ограничения $k$:
\begin{enumerate}
  \item Вычислить текущую невязку:
    \begin{equation}
      \delta\lambda_k = \frac{-(\mathbf{J}_k \mathbf{v} + b_k)}{A_{kk}}.
      \label{eq:pgs-delta}
    \end{equation}
  \item Обновить аккумулированный импульс с проекцией на допустимое множество:
    \begin{equation}
      \lambda_k \leftarrow \text{clamp}\!\left(\lambda_k + \delta\lambda_k,\; \lambda_k^{\min},\; \lambda_k^{\max}\right).
      \label{eq:pgs-clamp}
    \end{equation}
  \item Применить поправку скоростей:
    \begin{equation}
      \mathbf{v} \leftarrow \mathbf{v} + \mathbf{M}^{-1}\mathbf{J}_k^T \Delta\lambda_k,
      \label{eq:pgs-update}
    \end{equation}
    где $\Delta\lambda_k = (\lambda_k^{\text{new}} - \lambda_k^{\text{old}})$.
\end{enumerate}

Для нормального контакта: $[\lambda_k^{\min}, \lambda_k^{\max}] = [0, +\infty)$.
Для трения: $[-\mu\lambda_n, +\mu\lambda_n]$ (зависит от текущего нормального импульса).

Число итераций PGS является ключевым параметром производительности.
Обычно используют 10–20 итераций для интерактивных симуляций;
для точных сцен — 50–100 итераций.

\subsection{Warm Starting}

Для ускорения сходимости применяют технику \emph{warm starting}~\cite{catto2005iterative}:
импульсы $\boldsymbol{\lambda}$ от предыдущего шага сохраняются и применяются в начале нового шага как начальное приближение.
Это особенно эффективно при устойчивых контактах (раздел~2.2.4),
когда распределение импульсов меняется незначительно от шага к шагу.

\subsection{Стабилизация Баумгарте}

Из-за накопления численных ошибок тела могут «проникать» друг в друга.
Для устранения дрейфа вводят поправочный член Баумгарте в скоростное ограничение:
\begin{equation}
  b_k = \frac{\beta}{h}\,C_k,
  \label{eq:baumgarte}
\end{equation}
где $C_k$ — позиционная ошибка ограничения $k$, $\beta \in [0, 1]$ — коэффициент стабилизации.

При $\beta = 0$ коррекции позиции нет; $\beta = 1$ означает полное исправление за один шаг (может вызывать перерегулирование).
Типичное значение $\beta = 0.1$–$0.2$ обеспечивает плавную стабилизацию.

\subsection{Выводы по разделу}

В данном разделе изложена теоретическая основа импульсного решателя ограничений.
Контактная динамика сформулирована как задача LCP;
итерационный решатель PGS обеспечивает быстрое приближённое решение, достаточное для интерактивной симуляции.
Warm starting ускоряет сходимость для устойчивых контактов.

Совокупность методов, описанных в данной главе — симплектический интегратор Эйлера,
алгоритм GJK/EPA и PGS-решатель — образует стандартный конвейер физического движка реального времени.
Его реализация и адаптация для GPU-параллелизма составляет предмет следующих глав диссертации.
